\section{Experimental Setup}
\label{sec:experiments}

\subsection{Models and Baselines}
% Original placeholder retained: \todo{Describe the evaluated systems and justify the baseline selection.}
This section specifies a model-selection plan. No experiments have been completed for this draft. We propose a small set that separates visual understanding from native audiovisual understanding. Qwen2.5-VL-7B-Instruct~\citep{bai2025qwen25vl} and InternVL3.5-8B~\citep{wang2025internvl35} are candidate open visual baselines. Qwen2.5-Omni-7B~\citep{xu2025qwen25omni} is a candidate native audiovisual baseline. Their purpose is to provide reproducible reference systems with different input capabilities; this shortlist does not claim to cover the strongest available models.

We also propose one proprietary video-capable API, with an additional native audiovisual API only if needed to cover that input setting. \todo{Confirm the final models, exact checkpoint revisions, proprietary API snapshot identifiers, access conditions, and inference budget.} Model selection will be frozen before the main evaluation. A model without native audio input will not be presented as an audiovisual baseline simply because it receives a transcript.

\subsection{Input Conditions and Diagnostic Controls}
We propose modality-separated evaluation, following the general distinction used in v-HUB~\citep{shi2025vhub}. The \textbf{V} condition provides silent video frames. Text embedded in those frames remains visible. The \textbf{VA} condition provides frames and native audio to models that support both. The \textbf{T} condition provides supplied text, such as a controlled description or cue list, with no frames or audio. T is a diagnostic condition whose construction remains to be decided, rather than a substitute for audiovisual evaluation. \todo{Confirm the input conditions and the content and authorship of supplied text.}

Externally extracted speech transcripts or OCR outputs constitute separate \emph{ASR/OCR-assisted} conditions. Video titles, descriptions, and other source metadata would likewise belong to an explicitly identified metadata-assisted condition. These inputs must not be silently added to V or VA. Any comparison of V and VA should use the same clips and model where possible; comparisons between different models cannot isolate the effect of audio alone.

Reference labels should be condition-specific. A cue supported only by
audio should not count as an ordinary recognition miss in V. We propose
scoring claims supported by the supplied modalities and reporting excluded
or unanswerable items separately. An oracle cue may have an empty evidence
set after the supporting channel is removed.
\todo{Define modality-specific answerability, reference masks, and coverage reporting before evaluation.}

If the candidate relation task is retained, we propose both an
\emph{oracle-node} diagnostic, in which annotated target nodes are supplied,
and a candidate \emph{end-to-end} condition, in which the model identifies
instances before predicting relations. These answer different questions:
reasoning over known elements versus recovering the full structure from
video. The latter requires instance resolution beyond the concept-level
recognition task. A \emph{node-text-only} control will test how much relation
information is already exposed by supplied names or descriptions. Oracle
prompts should omit relation labels, role labels, and explanatory notes
that reveal the target. Even neutral descriptions may make roles
predictable, so this control is necessary before interpreting oracle
performance as video reasoning.
\todo{Confirm oracle input fields, instance-resolution outputs, the leakage control, and whether to retain the end-to-end relation condition.}

\subsection{Implementation Details}
% Original placeholder retained: \todo{Document prompts, decoding settings, model versions, and evaluation procedures needed for reproducibility.}
Implementation remains pending. \todo{Specify frame sampling and resolution, audio processing, context and output budgets, prompts, decoding settings, software versions, hardware, and cost.} These settings should be reported for each model and condition, including unsupported inputs and any truncation. Task metrics are proposed in Section~\ref{sec:tasks}; result tables and claims about model capabilities will be added only after evaluation.
